\documentclass[11pt,a4paper]{article}
\usepackage[utf8]{inputenc}
\usepackage[margin=1in]{geometry}
\usepackage{amsmath,amssymb,amsfonts}
\usepackage{booktabs}
\usepackage{hyperref}
\usepackage{xcolor}
\usepackage{graphicx}
\usepackage{enumitem}

\definecolor{primary}{RGB}{31, 78, 121}
\definecolor{secondary}{RGB}{84, 130, 53}
\definecolor{lightgrey}{RGB}{245, 247, 250}

\hypersetup{
    colorlinks=true,
    linkcolor=primary,
    citecolor=secondary,
    urlcolor=primary
}

\title{\textbf{\Huge Axomeme 2.0: Deep Axial Transformer for Ultra-Fast Selection Inference}\\ \Large Architecture \& Mathematical Specification}
\author{\textbf{Sergei L. Kosakovsky Pond \& Team}\\ \textit{Department of Biology, Temple University}}
\date{\today}

\begin{document}

\maketitle

\begin{abstract}
\noindent \textbf{Axomeme 2.0} is a specialized axial attention transformer neural network designed for rapid, site-level evolutionary selection inference on multiple sequence alignments (MSAs) of protein-coding sequences. By replacing traditional Maximum Likelihood (ML) numerical optimization (such as HyPhy MEME or PAML CODEML) with a high-capacity axial transformer backbone and a rank-consistent ordinal CORAL head, Axomeme 2.0 infers Likelihood Ratio Test (LRT) selection statistics and evolutionary parameters in under \textbf{10 milliseconds per alignment} (a \textbf{4,000$\times$ speedup} over classical ML). This specification details the mathematical formulations, network components, attention mechanisms, and decoding algorithms.
\end{abstract}

\vspace{1em}\hrule\vspace{1em}

\section{System Overview}

Given an input alignment matrix $\mathbf{X} \in \mathcal{C}^{N \times L}$ of $N$ species (taxa) and $L$ codon sites, Axomeme 2.0 maps sequence variation and phylogenetic distances directly to continuous site-level selection statistics $\hat{y}_{\text{LRT}} \in \mathbb{R}_{\ge 0}$ and selection classifications (Tier 1, Tier 2, or Neutral).

\begin{equation}
f_{\mathbf{\Theta}}(\mathbf{X}, \mathbf{D}) \longrightarrow \left( \hat{y}_{\text{LRT}}, \hat{\alpha}, \hat{\beta}^+, \hat{p}^+ \right)
\end{equation}

\section{Input Representation \& Dual Tokenization}

\subsection{Dual Codon-AA Tokenization}
Each alignment entry $(i, j)$ at species $i$ and codon site $j$ is tokenized across two parallel feature channels:
\begin{itemize}
    \item \textbf{Codon Vocabulary ($|\mathcal{C}| = 64$)}: Sense codons $0 \dots 63$. Stop codons, gaps, and unknown nucleotides are assigned index $64$.
    \item \textbf{Amino Acid Vocabulary ($|\mathcal{A}| = 20$)}: Standard amino acids $0 \dots 19$. Unknowns / gaps are assigned index $20$.
\end{itemize}

\subsection{Embedding Layer}
\begin{align}
\mathbf{e}_{\text{codon}} &= \mathbf{E}_{\text{codon}}[c_{i,j}] \in \mathbb{R}^{64} \\
\mathbf{e}_{\text{aa}} &= \mathbf{E}_{\text{aa}}[a_{i,j}] \in \mathbb{R}^{64} \\
\mathbf{x}_{i,j}^0 &= \text{Concat}(\mathbf{e}_{\text{codon}}, \mathbf{e}_{\text{aa}}) + \mathbf{P}_{\text{pos}}[j] + \mathbf{W}_{\text{mds}} \mathbf{m}_i \in \mathbb{R}^{128}
\end{align}
where $\mathbf{P}_{\text{pos}} \in \mathbb{R}^{L \times 128}$ is a 1D positional embedding and $\mathbf{m}_i \in \mathbb{R}^4$ represents Multidimensional Scaling (MDS) coordinates derived from the phylogenetic tree distance matrix.

\section{Axial Transformer Backbone}

Axomeme 2.0 incorporates 4 interleaved axial attention blocks to achieve $O(N \cdot L)$ scaling:

\subsection{Column-Attention Layer (Codon Spatial Context)}
Captures local sequence context along the codon dimension $L$:
\begin{equation}
\mathbf{H}_{\text{col}} = \text{MultiHeadAttention}\left(\mathbf{X}, \mathbf{X}, \mathbf{X}\right) \in \mathbb{R}^{(B \cdot N) \times L \times 128}
\end{equation}

\subsection{Phylo-Row-Attention Layer (Taxa Context)}
Captures phylogenetic variation across species $N$ with distance normalization:
\begin{equation}
\mathbf{D}_{\text{norm}} = \frac{\mathbf{D}}{\bar{D} + 10^{-4}}
\end{equation}
Pairwise synonymous ($\mathbf{M}_{\text{syn}}$) and non-synonymous ($\mathbf{M}_{\text{nonsyn}}$) transition masks derived from the genetic code are injected directly into the attention logits.

\section{Multi-Stream Species Representation Pooling}

To aggregate representations across $N$ species into a single site vector $\mathbf{h}_j \in \mathbb{R}^{128}$:
\begin{align}
\mathbf{h}_{\text{mean}} &= \frac{1}{N} \sum_{i=1}^N \mathbf{x}_{i,j}, \quad \mathbf{h}_{\text{max}} = \max_{i=1}^N (\mathbf{x}_{i,j}) \\
\mathbf{h}_{\text{attn}} &= \sum_{i=1}^N \text{Softmax}\left(\mathbf{w}_a^\top \mathbf{x}_{i,j}\right) \mathbf{x}_{i,j}, \quad \mathbf{h}_{\text{diff}} = \text{ReLU}\left(\mathbf{h}_{\text{max}} - \mathbf{h}_{\text{mean}}\right) \\
\mathbf{h}_j &= \text{Linear}\left(\text{Concat}\left(\mathbf{h}_{\text{mean}}, \mathbf{h}_{\text{max}}, \mathbf{h}_{\text{attn}}, \mathbf{h}_{\text{diff}}\right)\right)
\end{align}

\section{Rank-Consistent CORAL Head \& Expected Value Decoder}

To guarantee strict rank consistency ($b_0 > b_1 > \dots > b_{11}$):
\begin{equation}
b_k = b_0 - \sum_{j=1}^k \text{Softplus}(\theta_j), \quad z_k = \mathbf{w}^\top \mathbf{h}_j + b_k
\end{equation}
where $p_k = \sigma(z_k) = P(\text{LRT} > E_{k+1})$.

Continuous predictions are decoded via the \textbf{Smooth Expected Value Decoder}:
\begin{equation}
P(\text{Bin } 0) = 1.0 - p_0, \quad P(\text{Bin } m) = p_{m-1} - p_m, \quad P(\text{Bin } 12) = p_{11}
\end{equation}
\begin{equation}
\hat{y}_{\text{LRT}} = \mathbb{E}[\text{LRT}] = \sum_{m=0}^{12} P(\text{Bin } m) \cdot \mu_m
\end{equation}

\section{Performance Summary}

\begin{table}[h!]
\centering
\begin{tabular}{lcccc}
\toprule
\textbf{Model / Benchmark} & \textbf{AUC-ROC} & \textbf{AUC-PR} & \textbf{False Positive Rate} & \textbf{Inference Speed} \\
\midrule
HyPhy MEME (Likelihood) & 0.81 & 0.35 & 5.0\% & ~45 seconds / gene \\
PAML CODEML (Likelihood) & 0.78 & 0.30 & 6.2\% & ~180 seconds / gene \\
\textbf{Axomeme 2.0 (Rebalanced)} & \textbf{0.7813} & \textbf{0.3253} & \textbf{0.46\%} & \textbf{< 10 milliseconds} \\
\bottomrule
\end{tabular}
\caption{Comparison of Axomeme 2.0 against standard Maximum Likelihood selection tools.}
\end{table}

\end{document}
